\honest{Words as Hidden Inferences}{Eliezer Yudkowsky}{2008}

I reach into a barrel and pull out 11 blue eggs and 8 red cubes. When I feel another egg-shaped object, I guess ``blue,'' knowing that 19 is a small sample and that I am guessing. But when a tiger leaps out at me, I do not reason about the past properties of large striped feline shapes; I think ``Yikes! A tiger!'' The brain ``seems'' to have been adapted to make such inferences quickly and without tracking its assumptions. And once I name the objects ``bleggs'' and ``rubes,'' I stop noticing that I infer anything.

So it is ``a common misconception that you can define a word any way you like.'' That would be true only if words were logical classes from which you ``never took out any more information than you put in.'' But the brain categorizes whether or not we approve. \nb{The post cites no study. Research on infants supports it: Waxman and Markow (1995) found that new words direct twelve-month-olds' attention to categories of objects.}

The Aristotelians are my example. ``Aristotle knew perfectly well that Socrates was a human,'' though his own rules forbade concluding it before Socrates died. \nb{Socrates died fifteen years before Aristotle was born, so Aristotle never saw him. And Aristotle's own examples do not put mortality into the definition of man; they call man a two-footed or walking animal and list mortality as a property. The definition with ``mortal'' in it is Porphyry's, adopted by later logicians.} The Aristotelians went on recognizing people and bananas as everyone does, but had ``an erroneous picture of what they were doing'': asked how they knew Carol the grocer was mortal, they would cite the definition. \nb{Aristotle's own account in the \textsc{Posterior Analytics} is close to the post's: perception of a particular already presents ``man,'' and first premises are known by induction. He did hold that the resulting knowledge is always true.}

Mistaken theories of the mind do not interfere with quick perception, just as Aristotle did not die when he decided the brain was for cooling the blood. But they ``can severely mess up'' the deliberate thinking we use to correct first impressions. If you think you can define words as you like, you will not take the trouble to choose them wisely. \nb{The post gives no case of such harm; its only historical example is a mistaken theory that did none. Later posts in the sequence give examples.}
